\section{System and the Perception-to-Actuation Pipeline}
\label{sec:system}

\subsection{Vehicle and basin}

The physical platform is a BlueROV2 Heavy in the eight-thruster
configuration, running ArduSub over a tether, with a single forward
\num{1080}\,px camera streamed to the topside computer. It carries no
Doppler velocity log and no acoustic positioning, so it has no source of
absolute or velocity-referenced pose; this fact recurs throughout the
paper and eventually determines which of the two planners could be flown
at all. The basin is a rectangular test pool of roughly
\SI{6}{\metre}$\times$\SIrange{2.7}{3.0}{\metre}, with the obstacle -- a
\SI{0.5}{\metre}-class anchor -- suspended from a transverse rod.

\subsection{The pipeline and the boundary that both domains share}

The stack is a chain of ROS\,2 nodes, and the same source code runs in
both domains between two fixed points:
\[
\texttt{/perception/obstacles\_raw} \rightarrow \text{qualification}
\rightarrow \text{estimator} \rightarrow \text{planner} \rightarrow
\texttt{/planner/cmd\_vel\_safe}.
\]
Everything inside that boundary is identical, file for file, in
simulation and in the water. What differs by construction is what
produces the raw observations at one end and what consumes the safe
command at the other. Fixing the boundary is what makes a simulated and a
physical run comparable at all; Section~\ref{sec:deployment} reports the
seven ways in which, in the event, the two sides ceased to match.

\emph{Qualification} rejects observations that are not yet trustworthy.
It requires a number of mutually coherent updates before declaring a
track valid, which suppresses single-frame false positives at the cost of
a delay before the planner may act.

\emph{Estimation} is a fixed-gain Kalman filter on the obstacle's image
position and extent, which also predicts through gaps in the observation
stream. Earlier work in this project compared four estimator variants;
the results are summarised in Section~\ref{sec:frozen} but are not the
subject of this paper, and the estimator is held fixed at the same
variant everywhere.

\emph{Range} is monocular and deliberately simple. An object of assumed
physical height $H$ whose bounding box fills a fraction $h$ of the image
subtends $\theta = h \cdot \mathrm{vfov}$, giving
\begin{equation}
R = \frac{H/2}{\tan(\theta/2)} .
\label{eq:range}
\end{equation}
Both planners call the same implementation of
Equation~\ref{eq:range} with the same assumed height, so any range error
is shared and cannot advantage either. We note here, and return to it in
Section~\ref{sec:limitations}, that the vertical field of view used in
Equation~\ref{eq:range} was \SI{90}{\degree} in both domains -- the node
default -- and not the \SI{60}{\degree} recorded in our own frozen
benchmark file, which no node reads.

\subsection{The two planners}

The \emph{committed} planner is a finite-state controller. Below an
engagement range it commits to one side, holds that side for a minimum
duration, strafes to a clearing offset, runs past the obstacle and
returns to the original line. Its side choice depends only on where the
obstacle appears in the image, which is the property
Section~\ref{sec:realresults} tests. Its manoeuvre geometry -- the
clearing offset and the run-past margin -- is the pair of parameters that
turns out not to fit the basin.

The \emph{holonomic dynamic-window} planner is the standard local
baseline~\cite{fox1997thedynam}, adapted to a holonomic vehicle and
given the obstacle radius and goal lookahead appropriate to the pool. It
scores sampled velocity candidates for clearance, progress, speed, route
adherence and smoothness. Crucially for what follows, it requires a pose
estimate to evaluate route adherence and to project candidates forward.

\subsection{Actuation}

Simulated commands pass to a vehicle-dynamics model in HoloOcean.
Physical commands are mapped to ArduSub joystick channels through an
affine deadband model,
\begin{equation}
c = \operatorname{sign}(v)\left(d + \frac{|v|}{k}\right),
\label{eq:deadband}
\end{equation}
with per-direction gain $k$ and deadband $d$ identified from open-water
trials (Section~\ref{sec:calibration}). The same identified model is what
the S3 calibration level injects into simulation, downstream of the
planner, so that a calibration level never edits the controller.
